% 原书第 6—11 页；第 1.3—1.4 节。
\section{靶向学习的实际应用}
\label{sec:ch01-in-practice}

在各种现实分析中，已有许多广受关注的例子展示了靶向学习的益处。例如，Pandora 的数据科学家采用靶向学习发现，流媒体音乐的播放使新发行音乐的销量提高了 2.3\%，使既有曲目销量提高了 2.7\%。《Billboard》杂志曾报道这项研究。医疗保健是应用靶向学习尤为频繁的一个领域。靶向学习如何改善医疗保健？在发表于《Lancet Respiratory Medicine》的研究中，研究人员为重症监护病房开发了一种预测死亡率的超级学习器，其表现优于疾病严重程度评分。临床医生可通过在线用户界面使用这一算法。在另一项发表于《World Psychiatry》的研究中，研究人员构建了一个新函数，用于预测创伤事件后的创伤后应激障碍（posttraumatic stress disorder，PTSD）。该算法的表现十分出色，96\% 的 PTSD 结局都落在预测值最高的 10\% 人群中。近期，一项计算卫生经济学分析利用大型健康记录理赔数据库，考察了单项病症对医疗总支出的影响。用于效应估计的靶向学习估计量，在控制人口学信息和其他病症后，按照各病症类别对医疗总支出的贡献对其排序。以往关于病症对医疗支出影响的研究，主要采用医保计划支付风险调整中的参数回归公式，或未经混杂因素调整的汇总均值。这项研究的结果表明，多发性硬化、充血性心力衰竭、重症癌症、重性抑郁和双相障碍，以及慢性肝炎，是费用最高的病症（见表 1.1）。相比之下，用于医保计划支付风险调整的参数回归公式，无论在效应估计值的大小，还是在相对排序上，都存在不可忽视的差异。如果现行风险调整公式未能准确反映各病症的增量效应，医疗保险机构挑选参保人的激励就可能继续存在。关于早期研究的其他参考文献，见第 1.6 节。

\begin{table}[htbp]
 \centering
 \caption{靶向学习效应估计值最高的五类病症及其参数回归估计值}
 \label{tab:ch01-cost-estimates}
 \begin{tabular}{lrr}
  \toprule
  病症 & 靶向学习 & 参数回归\\
  \midrule
  多发性硬化 & 67,011 & 30,715\\
  充血性心力衰竭 & 19,904 & 4131\\
  肺癌、脑癌及重症癌症 & 19,233 & 24,528\\
  重性抑郁与双相障碍 & 15,398 & 3498\\
  慢性肝炎 & 10,530 & 5539\\
  \bottomrule
 \end{tabular}
\end{table}

\section{统计估计问题}
\label{sec:ch01-statistical-estimation}

我们用一个简化的体外受精（in vitro fertilization，IVF）例子，介绍纵向统计估计问题。这里，我们在控制时变混杂因素（time-dependent confounder）的情况下，估计一项最多包含两个 IVF 周期的治疗计划获得成功的概率，即胚胎移植后活产的概率。不孕症是一个全球性的公共卫生问题，目前已有多种治疗方法可供选择。IVF 是一种日益普遍的治疗方法，但由于 IVF 治疗计划包含多个周期，要准确评价其成功率一直颇具挑战。

\subsection{数据}
\label{subsec:ch01-data}

考虑各时点 $t$（$t=0,\ldots,T+1$）的协变量（covariate）向量 $L_t$，其中，基线协变量记为 $L_0$。为了说明问题，我们将 IVF 研究限定为一种特定的数据结构，其中只有两个相继的时点接受干预。这一数据结构是横断面数据的简单扩展；横断面情形只在单个时点接受干预。

在这一数据结构中，各时点 $t$ 的协变量组成 $\bar L=(L_0,L_1,L_2)$，且 $T=1$。在序贯决策过程文献中，协变量集合 $\bar L$ 也称为状态集合；此外，也常使用 $L_{1:2}$ 之类的下标记号来表示协变量集合的某个特定子集。$L_0$ 中的协变量包括母亲年龄、IVF 机构、获卵数、移植或冷冻的胚胎数，以及妊娠指标。$L_1$ 和 $L_2$ 编码 IVF 周期是否成功。干预集合（也称为“行动”）记为 $\bar A=(A_0,A_1)$，其中随机变量 $A_t$ 表示时点 $t$ 的干预。在 IVF 研究中，$\bar A$ 中的每个 $A_t$ 都是二值变量，表示是否尝试了该 IVF 周期。按约定，如果 $A_t=0$，则 $L_{t+1}=0$；如果 $L_1=1$，则 $L_2=1$。

于是，我们可以把一名随机抽取的研究对象的数据表示为按时间排序的数据结构：
\[
 O=(L_0,A_0,L_1,A_1,L_2),
\]
其中，假定 $L_t$ 发生在 $A_t$ 之前。最终的测量值 $L_2$ 记为 $Y$，表示感兴趣的结局。为简便起见，我们考虑 $Y$ 为二值变量的情形。样本由随机变量 $O$ 的 $n$ 次独立同分布（independent and identically distributed，i.i.d.）抽样 $O_1,\ldots,O_n$ 构成。这些随机变量的实现值记为 $o_1,\ldots,o_n$。按照数据的时间顺序，$O$ 的概率分布可分解为
\[
 \begin{aligned}
 p(O)={}&p(L_0)\times p(L_1\mid A_0,L_0)
                \times p(Y\mid A_1,L_1,A_0,L_0)\\
        &{}\times p(A_0\mid L_0)\times p(A_1\mid L_1,A_0,L_0).
 \end{aligned}
\]

\subsection{模型与参数}
\label{subsec:ch01-model-parameter}

我们假设，观测数据的生成分布属于一个非参数统计模型 $\mathcal M$；这一模型包含可能的概率分布集合。感兴趣的参数，是一项最多包含两个 IVF 周期的治疗计划获得成功的概率，即胚胎移植后活产的概率，它可以写成
% 此处严格照录印刷原式：求和号无指标，第二因子为 L_1=1。
\[
 \begin{aligned}
 \Psi(P_0)=E_{P_0}\Bigl(&\sum P_0(Y=1\mid A_{0:1}=1,L_{0:1}=l_{0:1})\\
 &{}\times P_0(L_1=1\mid A_0=1,L_0=l_0)\Bigr).
 \end{aligned}
\]
在第 2 章将讨论的因果可识别性假设下，因果参数可以写为
\[
 P(Y_{(1,1)}=1),
\]
且等于 $\Psi(P_0)$，其中 $Y_{(1,1)}$ 是干预 $A_{0:1}=1$ 下的反事实结局（counterfactual outcome）。

\subsection{靶向最小损失估计量}
\label{subsec:ch01-tmle}

针对这一研究问题，可以构造靶向最小损失估计量（targeted minimum loss-based estimator，TMLE），以估计 $P(Y_{(1,1)}=1)$。TMLE 框架是一个极其一般的系统，由损失函数（loss function）、初始估计量，以及通过初始（或当前）估计量的最不利子模型（least favorable submodel）所定义。确切地说，它需要以下要素：
\begin{enumerate}
 \item 定义一个目标参数：它将观测单位数据结构的概率分布的一个参数（通常为无限维参数）映射到参数空间。
 \item 推导目标参数路径导数的有效影响曲线。
 \item 规定一个损失函数。
 \item 指定一个通过该参数的初始（或当前）估计量的最不利子模型，使扰动（fluctuation）为零时，基于损失的得分的线性张成空间包含有效影响曲线。
 \item 给出一个算法：对最不利参数子模型的扰动参数，迭代地最小化由指定损失函数定义的经验风险（empirical risk），并更新初始（或当前）估计量。
\end{enumerate}

这一迭代最小化过程一直进行到扰动参数的极大似然估计值接近零。根据子模型和损失函数所满足的广义基于损失的得分条件，所得无限维参数的 TMLE 满足有效得分方程（efficient score equation）。这为目标参数相应的代入估计量具有双重稳健性和渐近有效性（asymptotic efficiency）奠定了基础。靶向极大似然估计量是 TMLE 的一种。第 5 章将证明，总能选取一个最不利子模型，使得仅需对由指定损失函数定义的经验风险进行一次最小化，就足以求解有效得分方程。这些特殊的局部最不利子模型称为通用最不利子模型（universal least favorable submodel）。

\subsection{其他常见的估计问题}
\label{subsec:ch01-other-problems}

\noindent\textbf{特定规则下的均值（rule-specific mean）。}
在利用纵向观察数据研究临床问题时，我们常常希望评价随时间延续的处理规则，这类规则称为动态处理规则（dynamic rule），有时也称为处理策略或处理方案（regime 或 regimen）。从更广泛的数理科学视角来看，动态处理规则属于序贯决策理论。其中的一个关键点是：对于某个给定时点，如果考虑即时结局和考虑长期结局，最优的干预决策可能不同。这类决策在统计学中十分常见，尤其常见于医学和公共卫生问题，因此，估计最优动态处理规则的方法具有重要意义。所考虑的规则还可能包含复杂的分类器。例如，考虑美国既往未接受治疗的 HIV 感染者应在何时开始抗逆转录病毒治疗的问题。在这里，我们希望考虑一组预先指定的 CD4 计数阈值，例如从 $200$ 到 $500\ \text{细胞}/\mathrm{mm}^3$，每隔 $50$ 设置一个阈值；当 CD4 计数低于阈值时，就应开始抗逆转录病毒治疗。每个动态处理规则都以一个 CD4 阈值作为索引。

设 $\mathcal D=(d_1,\ldots,d_K)$ 是我们考虑的动态处理规则集合。每个动态处理规则 $d_k$ 编码一个按时间排列的规则序列 $d_{k,t}$（$t=1,\ldots,T$），其中，$d_{k,t}$ 表示将患者此前测得的协变量 $\bar L_t$ 映射为时点 $t$ 应接受的处理 $a(t)$ 的函数。在不会引起歧义时，为简化记号，我们省略单个规则 $d$ 的下标。如果某个个体实际接受的干预与规则 $d$ 指示的干预相同，就称该个体直到时点 $t$ 都遵循规则 $d$。设 $Y^d$ 表示假如令该研究对象在所有时点都遵循规则 $d$，本会观测到的（反事实）取值。我们的目标是确定最优动态处理规则 $d^*$，即 $\mathcal D$ 中使潜在结局（potential outcome）$Y^d$ 的期望值最大的规则。我们首先考虑这样一个估计问题：对每个规则 $d\in\mathcal D$，估计假如所有人都遵循规则 $d$ 时，$Y$ 的总体均值，即 $E(Y^d)$。由于本例数据结构中的 $Y$ 是二值变量，这也等于 $Y^d=1$ 的概率。在一组较强的可识别性假设下，包括无未测量混杂因素的假设，可以利用 g 计算公式（g-computation formula），将 $E(Y^d)$ 表示为观测数据生成分布 $P$ 的一个函数：
\[
 \begin{aligned}
 \Psi^d(P)={}&\sum_{l_0}\sum_{l_1}
 P\bigl(Y=1\mid A_1=d(\bar l_1),\bar L_1=\bar l_1,A_0=d(l_0)\bigr)\\
 &{}\times P\bigl(L_1=l_1\mid A_0=d(l_0),L_0=l_0\bigr)P(L_0=l_0).
 \end{aligned}
\]
我们在这里较为详细地说明动态处理规则的数据结构和统计估计问题，是因为第 2—4 章的多处内容将把这项研究作为贯穿的实例来讲解。

\medskip
\noindent\textbf{社区随机试验。}
在社区随机试验（community randomized trial）中，我们关心的是这样一种情形：将干预分配给少数几个社区，并从每个社区随机抽取个体样本，收集其协变量和结局数据。借助我们此前关于有偏抽样的研究，可以在这些试验中识别感兴趣的参数，例如因果对比。第 13 章将更详细地介绍这一数据结构。

\medskip
\noindent\textbf{网络。}
网络数据日益常见。从最初的疾病传播，到后来的电子邮件，再到如今迅速扩张的社交网络，我们经常观测到这样的总体：其中各个单位按照某种网络，在因果关系上相互连接。其数据结构通常是纵向的，包含时变暴露和协变量。第 20 章和第 21 章将更详细地介绍这些数据结构和参数。
